\documentclass[11pt,a4paper]{article}
\usepackage[margin=1.85cm]{geometry}
\usepackage{amsmath,amssymb,mathtools,booktabs,array,microtype}
\usepackage{xcolor}
\usepackage{enumitem}
\usepackage[round,authoryear]{natbib}
\usepackage[colorlinks=true,linkcolor=blue!45!black,citecolor=blue!45!black,urlcolor=blue!45!black]{hyperref}
\setlist{nosep,leftmargin=*}
\allowdisplaybreaks

\newcommand{\kb}{k_{\mathrm B}}
\newcommand{\K}{\mathcal K}
\newcommand{\half}{\tfrac12}
\newcommand{\code}[1]{\texttt{#1}}
\newcommand{\pos}[1]{\left[#1\right]_{+}}
\newcommand{\dd}{\mathrm d}
\newcommand{\e}{\mathrm e}
\newcommand{\rc}{\textsf{RadCluster}}
\definecolor{warn}{RGB}{252,244,204}

\title{\vspace{-1.2cm}\textbf{Systematic Derivation of the Three
$\half\langle111\rangle\!\rightarrow\!\langle100\rangle$ Loop-Conversion Kernels}\\[3pt]
\large Thermodynamic rotation, collision-induced junction formation, and absorption growth}
\author{Technical derivation for the \rc{} formulation}
\date{15 September 2026}

\begin{document}
\maketitle
\vspace{-1em}

\begin{abstract}
This note derives, step by step, the three rate kernels used to transfer
self-interstitial content between mobile $\half a\langle111\rangle$ loops and
sessile $a\langle100\rangle$ loops in \rc.  The unary kernel is constructed by
combining the anisotropic-elastic free-energy difference of
\citet{DudarevBullough2008} with transition-state kinetics and a one-way net-rate
closure.  The junction kernel is obtained by multiplying the binary encounter
kernel by (i) a size-and-orientation compatibility factor and (ii) the
probability that the metastable $\half\langle110\rangle$ intermediate of
\citet{MarianWirth2002} reaches $\langle100\rangle$ before reverting.  The
absorption kernel follows from the same diffusion-limited encounter problem,
but only the $\half\langle111\rangle$ partner is mobile and the pre-existing
$\langle100\rangle$ loop templates the product character.  Every empirical or
calibrated closure is distinguished from quantities taken from elasticity,
atomistics, or geometry.  Numerical checks expose a parameter inconsistency in
the supplied draft and give the corrected barrier law.
\end{abstract}

\section{Reaction network, state variables, and units}
\label{sec:setup}

Let $c_n^{111}$ and $c_m^{100}$ denote the atomic fractions (number per lattice
site) of interstitial clusters containing $n$ and $m$ SIAs.  If number densities
are used instead, write $C_n=c_n/\Omega$, where $\Omega$ is the atomic volume.
The three admissible reactions are
\begin{align}
  111_n &\longrightarrow 100_n,                                      &&\text{unary rotation}, \label{rxn:uni}\\
  111_n+111_{n'} &\longrightarrow 100_{n+n'},                         &&\text{junction},       \label{rxn:junc}\\
  100_m+111_n &\longrightarrow 100_{m+n},                             &&\text{absorption}.     \label{rxn:abs}
\end{align}
Here $111_n$ abbreviates a $\half a\langle111\rangle$ loop and $100_m$ an
$a\langle100\rangle$ loop.  The $100$ population is defined only for
$m\ge n_{\rm loop}^{\min}$ and is sessile, $D_m^{100}=0$.  The $111$ population
is mobile only for $n\le i_{\rm mobile}$.

For atomic-fraction variables, a unary kernel has units s$^{-1}$ and a binary
kernel also has units s$^{-1}$:
\begin{equation}
  \dot c_a=-\Gamma_a c_a,qquad
  \dot c_a=-K_{ab}c_ac_b.                                             \label{eq:units}
\end{equation}
The corresponding volumetric binary coefficient is
$k_{ab}^{(V)}=\Omega K_{ab}$, with units m$^3$~s$^{-1}$, because
$\dot C_a=-k_{ab}^{(V)}C_aC_b$.

\begin{table}[ht]
\centering\small
\caption{Logical factorization of the three kernels.  Factors labeled
``closure'' are not uniquely determined by the cited microscopic theory.}
\label{tab:factorization}
\begin{tabular}{p{2.5cm}p{4.3cm}p{4.3cm}p{3.7cm}}
\toprule
Kernel & Attempt or encounter & Selection/success & Admissibility \\
\midrule
$\Gamma_{\rm uni}(n,T)$ & transition-state attempt rate & Dudarev driving-force gate (closure) & loop floor; $\Delta F>0$ \\
$K_{nn'}^{\rm junc}$ & $111$--$111$ encounter kernel & size/orientation compatibility and Marian completion probability (closures) & both partners mobile; size floor \\
$K_{mn}^{\rm abs}$ & mobile $111$ encounter with sessile $100$ & templated completion probability (closure) & $m\ge n_{\rm loop}^{\min}$; $n\le i_{\rm mobile}$; $m+n\le I$ \\
\bottomrule
\end{tabular}
\end{table}

\section{Common geometry and loop free energy}
\label{sec:geometry}

\subsection{Area, perimeter, and size scaling}

Insertion of $n$ atomic volumes produces a platelet volume $A_Xb_X=n\Omega$.
Thus, for character $X\in\{111,100\}$,
\begin{equation}
 A_X(n)=\frac{n\Omega}{b_X},\qquad
 b_{111}=\frac{\sqrt3}{2}a,\qquad b_{100}=a.                           \label{eq:area}
\end{equation}
For the hexagonal $111$ and square $100$ shapes used in \rc,
\begin{equation}
 P_{111}(n)=6\sqrt{\frac{2A_{111}(n)}{3\sqrt3}},\qquad
 P_{100}(n)=4\sqrt{A_{100}(n)}.                                      \label{eq:perimeters}
\end{equation}
Substitution of Eq.~\eqref{eq:area} into Eq.~\eqref{eq:perimeters} gives
\begin{align}
 \frac{P_{111}(n)}{b_{111}}
 &=\underbrace{\frac{6}{b_{111}}
 \sqrt{\frac{2\Omega}{3\sqrt3\,b_{111}}}}_{C_P}\sqrt n,
                                                               \label{eq:Cp}\\
 \frac{P_{100}(n)}{b_{100}}
 &=\frac{4}{b_{100}}\sqrt{\frac{\Omega}{b_{100}}}\sqrt n.     \label{eq:P100scale}
\end{align}
Hence all perimeter-proportional energetic terms scale as $\sqrt n$.

\subsection{Dudarev anisotropic-elastic free energy}

The loop formation free energy is represented by
\begin{equation}
 E_l^X(n,T)=P_X(n)\left\{
 \widehat F_X(T)\ln\!\left[\frac{4R_X^*(n)}{\e\delta}\right]
 +F_\delta^X+F_c^X\right\},                                         \label{eq:loopF}
\end{equation}
where $R_X^*=\sqrt{A_X/\pi}$ is an equal-area radius, $\delta$ is the core
cutoff, $\widehat F_X$ is the anisotropic prelogarithmic line-energy factor,
and $F_\delta^X+F_c^X$ represents core traction plus nonlinear core energy
\citep{DudarevBullough2008}.  The conversion driving force at fixed SIA content
is
\begin{equation}
 \Delta F(n,T)=E_l^{111}(n,T)-E_l^{100}(n,T).                          \label{eq:dF}
\end{equation}
The sign convention is important: $\Delta F>0$ means that the $100$ product is
thermodynamically lower in free energy.

For a pure-edge $001[100]$ dislocation, anisotropic elasticity gives
\begin{equation}
 \widehat F_{001}([100],T)=\frac{a^2}{4\pi}(c_{11}+c_{12})
 \left[\frac{c_{44}(c_{11}-c_{12})}
 {c_{11}(c_{11}+c_{12}+2c_{44})}\right]^{1/2}.                        \label{eq:F100anis}
\end{equation}
Near the bcc--fcc transformation, the tetragonal shear modulus obeys the
empirical softening law
\begin{equation}
 c'(T)=\frac{c_{11}-c_{12}}2\simeq c'_0(1-T/T_c)^{1/2},                \label{eq:cprime}
\end{equation}
so Eq.~\eqref{eq:F100anis} implies
\begin{equation}
 \widehat F_{100}(T)=\widehat F_{100}^{,0}
 (1-T/T_c)^{1/4}.                                                     \label{eq:F100T}
\end{equation}
This is the physical source of the strong temperature dependence: the $100$
elastic line energy softens, while the relevant $111$ factors remain finite
\citep{DudarevBullough2008}.

The fitted zero-temperature coefficient $\widehat F_{100}^{,0}$ is obtained
from one crossover condition, not independently prescribed.  Define
\begin{equation}
 L_X(n)=\ln\!\left[\frac{4R_X^*(n)}{\e\delta}\right],\qquad
 Q(T)=(1-T/T_c)^{1/4}.                                                \label{eq:LQ}
\end{equation}
Imposing $\Delta F(n_{\rm ref},T^*)=0$ in Eq.~\eqref{eq:dF} yields
\begin{equation}
 \boxed{
 \widehat F_{100}^{,0}=
 \frac{\dfrac{P_{111}}{P_{100}}
 \left[\widehat F_{111}L_{111}+F_\delta^{111}+F_c^{111}\right]
 -(F_\delta^{100}+F_c^{100})}
 {Q(T^*)L_{100}} }
 \quad\text{at }n=n_{\rm ref}.                                      \label{eq:F100cal}
\end{equation}
Equation~\eqref{eq:F100cal} documents exactly which datum is calibrated and
prevents $\widehat F_{100}^{,0}$ and $T^*$ from being counted as independent
inputs.

\section{Kernel 1: unary $111_n\rightarrow100_n$ rotation}
\label{sec:unary}

\subsection{Step 1: transition-state attempt rate}

For a $111$ loop trapped in a local free-energy minimum, transition-state theory
gives
\begin{equation}
 k_+(n,T)=\nu_0\exp\!\left[-\frac{E_a(n)}{\kb T}\right],              \label{eq:kplus}
\end{equation}
where $E_a=F^\ddagger-F_{111}$ is the forward activation barrier and $\nu_0$
is an attempt frequency.  The available atomistic literature motivates a
thermally activated reorientation pathway, but does not supply a unique barrier
for every loop size.  The \rc{} closure assumes that the activated portion of
the loop line grows with perimeter:
\begin{equation}
 E_a(n)=E_{a0}+\gamma_a\frac{P_{111}(n)}{b_{111}}.                     \label{eq:Ea}
\end{equation}
Here $P_{111}/b_{111}$ is the number of Burgers-vector-length line segments and
$\gamma_a$ is an effective barrier per segment.

\subsection{Step 2: thermodynamic direction and the net-rate factor}

For an explicitly reversible pair of states, local detailed balance requires
\begin{equation}
 \frac{k_+(n,T)}{k_-(n,T)}=
 \exp\!\left[\frac{\Delta F(n,T)}{\kb T}\right].                     \label{eq:DB}
\end{equation}
Choosing $k_-=k_+\exp[-\Delta F/(\kb T)]$ gives the forward-minus-reverse
coefficient
\begin{equation}
 k_+-k_-=k_+\left[1-\exp\!\left(-\frac{\Delta F}{\kb T}\right)\right]. \label{eq:net}
\end{equation}
\rc{} uses Eq.~\eqref{eq:net} as an effective irreversible rate only where the
forward reaction is downhill.  Defining $[x]_+=\max(0,x)$ therefore gives
\begin{equation}
 \boxed{
 \Gamma_{\rm uni}(n,T)=\nu_0\exp\!\left[-\frac{E_a(n)}{\kb T}\right]
 \pos{1-\exp\!\left[-\frac{\Delta F(n,T)}{\kb T}\right]}
 \Theta(n-n_{\rm loop}^{\min}) .}                                    \label{eq:Gamma}
\end{equation}
The bracket has the correct limiting behavior:
\begin{equation}
 \pos{1-\e^{-\Delta F/\kb T}}\simeq
 \begin{cases}
   0, & \Delta F\le0,\\
   \Delta F/(\kb T), & 0<\Delta F\ll\kb T,\\
   1, & \Delta F\gg\kb T.
 \end{cases}                                                         \label{eq:gateLimits}
\end{equation}
This construction is best described as a \emph{one-way net-rate closure
consistent with the sign of the Dudarev driving force}.  It is not a fully
reversible detailed-balance model, because the actual term
$+k_-c_n^{100}$ is not retained.  A reversible alternative would be
\begin{equation}
 \dot c_n^{111}=-k_+c_n^{111}+k_-c_n^{100},\qquad
 \dot c_n^{100}=+k_+c_n^{111}-k_-c_n^{100}.                            \label{eq:reversible}
\end{equation}

\subsection{Step 3: population equation and conservation}

The implemented unary contribution is
\begin{equation}
 \left.\dot c_n^{111}\right|_{\rm uni}=-\Gamma_{\rm uni}c_n^{111},
 \qquad
 \left.\dot c_n^{100}\right|_{\rm uni}=+\Gamma_{\rm uni}c_n^{111}.   \label{eq:uniME}
\end{equation}
Therefore
\begin{equation}
 \left.\frac{\dd}{\dd t}\sum_n n(c_n^{111}+c_n^{100})\right|_{\rm uni}=0, \label{eq:uniCons}
\end{equation}
and loop number is also conserved by this channel.

\subsection{Numerical reduction of the barrier}

Using $a=2.867$~\AA, $\Omega=11.80$~\AA$^3$, and Eq.~\eqref{eq:Cp} gives
\begin{equation}
 b_{111}=2.483~\text{\AA},\qquad C_P=3.26835,qquad
 \frac{P_{111}(n)}{b_{111}}=3.26835\sqrt n.                            \label{eq:CpNum}
\end{equation}
With the supplied table values $E_{a0}=2.10$~eV and
$\gamma_a=0.020$~eV,
\begin{equation}
 \boxed{E_a(n)=2.10+0.06537\sqrt n\ \text{eV}.}                       \label{eq:EaNum}
\end{equation}

\noindent\colorbox{warn}{\parbox{0.96\linewidth}{
\textbf{Consistency correction.}  The expression
$E_a=1.80+0.098\sqrt n$~eV in the supplied prose corresponds to the earlier
defaults $E_{a0}=1.8$~eV and $\gamma_a=0.03$~eV found in the current repository,
not to the supplied parameter table.  With that table, Eq.~\eqref{eq:EaNum} is
the correct reduction.  Similarly, if $\Delta F>0$ only for $1\le n\le23$ and
$n_{\rm loop}^{\min}=4$, the actual support of Eq.~\eqref{eq:Gamma} is
$4\le n\le23$.}}

\section{Kernel 2: $111_n+111_{n'}\rightarrow100_{n+n'}$ junction}
\label{sec:junction}

\subsection{Step 1: diffusion-limited encounter rate}

For two three-dimensionally diffusing objects, introduce the relative
coordinate ${\bf r}={\bf r}_1-{\bf r}_2$ and relative diffusivity
$D_{\rm rel}=D_n+D_{n'}$.  Outside an absorbing encounter surface of radius
$R_c=R_n^c+R_{n'}^c$, the steady pair concentration obeys
\begin{equation}
 D_{\rm rel}\nabla^2 C=0,\qquad C(R_c)=0,\qquad C(r\rightarrow\infty)=C_\infty. \label{eq:Laplace}
\end{equation}
The spherical solution is $C(r)=C_\infty(1-R_c/r)$.  Integrating the inward
flux at $R_c$ gives the Smoluchowski coefficient
\begin{equation}
 k_{nn'}^{(V)}=\int_{S_c}D_{\rm rel}\nabla C\cdot\dd{\bf S}/C_\infty
 =4\pi R_c(D_n+D_{n'}).                                                \label{eq:Smol}
\end{equation}
To reproduce the radius convention used in \rc{}, write
\begin{equation}
 R_n^c=2\Omega^{1/3}\xi_n,\qquad
 \xi_n=\left(\frac{3n}{8\pi}\right)^{1/3}.                            \label{eq:xi}
\end{equation}
Dividing Eq.~\eqref{eq:Smol} by $\Omega$ converts the volumetric coefficient to
the atomic-fraction kernel of Eq.~\eqref{eq:units}:
\begin{equation}
 \boxed{
 \K^{ii}_{n,n'}=\frac{8\pi}{\Omega^{2/3}}
 (\xi_n+\xi_{n'})(D_n^{111}+D_{n'}^{111}).}                            \label{eq:Kii}
\end{equation}
If $D^{111}$ is the isotropic-equivalent diffusivity of a 1D/3D glider, then
Eq.~\eqref{eq:Kii} is an effective-medium closure.  Pure 1D and mixed-dimensional
encounters generally require different first-passage coefficients; analytical
and kinetic Monte Carlo treatments are discussed by \citet{HuangMarian2019}
and the effect on ferritic-iron cluster dynamics by \citet{KohnertWirth2015}.

\subsection{Step 2: branching at encounter}

The Burgers-vector reaction identified by \citet{MarianWirth2002} is
\begin{equation}
 \half[111]+\half[1\bar1\bar1]\longrightarrow[100].                  \label{eq:bvec}
\end{equation}
Not every encounter has the required Burgers-vector combination, geometry,
impact parameter, or relative size.  The binary collision kernel is therefore
split without changing the total encounter budget:
\begin{align}
 K_{nn'}^{111\rightarrow100}&=\varphi_{\rm junc}(n,n',T)\K^{ii}_{n,n'}, \label{eq:Kjunc}\\
 K_{nn'}^{111\rightarrow111}&=[1-\varphi_{\rm junc}(n,n',T)]\K^{ii}_{n,n'}, \label{eq:Kself}\\
 K_{nn'}^{111\rightarrow100}+K_{nn'}^{111\rightarrow111}&=\K^{ii}_{n,n'}. \label{eq:budget}
\end{align}

\subsection{Step 3: completion probability of the metastable intermediate}

The Marian pathway contains a metastable $\half\langle110\rangle$ configuration.
From that state, let completion and reversal be independent activated events:
\begin{equation}
 k_2=\nu_2\e^{-\Delta H_2/\kb T},\qquad
 k_{\rm rev}=\nu_r\e^{-\Delta H_{\rm rev}/\kb T}.                     \label{eq:compRates}
\end{equation}
For two independent exponential waiting times, the probability that completion
occurs first is
\begin{align}
 P_{\rm succ}(T)
 &=\int_0^\infty\underbrace{k_2\e^{-k_2t}}_{\substack{\text{completion}\\
 \text{at }t}}\underbrace{\e^{-k_{\rm rev}t}}_{\substack{\text{no reversal}\\
 \text{before }t}}\dd t                                               \label{eq:hazardInt}\\
 &=\frac{k_2}{k_2+k_{\rm rev}}                                        \label{eq:hazard}\\
 &=\left[1+\frac{\nu_r}{\nu_2}
 \exp\!\left(\frac{\Delta H_2-\Delta H_{\rm rev}}{\kb T}\right)\right]^{-1}. \label{eq:Pgeneral}
\end{align}
Assuming equal attempt frequencies, $\nu_2=\nu_r$, gives the implemented gate
\begin{equation}
 \boxed{P_{\rm succ}(T)=
 \left[1+\exp\!\left(\frac{\Delta H_2-\Delta H_{\rm rev}}{\kb T}\right)\right]^{-1}.} \label{eq:Psucc}
\end{equation}
Only the barrier difference is identifiable from this scalar gate.  Separate
barriers would become identifiable only if an absolute residence time or total
escape rate from the intermediate were also used.

\subsection{Step 4: size compatibility and final junction kernel}

The unresolved geometrical compatibility is represented by a log-size Gaussian:
\begin{equation}
 g(n,n')=\exp\!\left[-\frac{\ln^2(n/n')}{2\sigma_s^2}\right].          \label{eq:gsize}
\end{equation}
It is symmetric, $g(n,n')=g(n',n)$; it peaks at equal size; and
$g=\e^{-1/2}$ at $n/n'=\e^{\pm\sigma_s}$.  Combining compatibility,
completion, and the minimum size gives
\begin{equation}
 \boxed{
 \varphi_{\rm junc}=P_{\rm succ}(T)\,\varphi_{\max}\,g(n,n')
 \Theta[\min(n,n')-n_{j}^{\min,\rm eff}],}                            \label{eq:phi}
\end{equation}
where
\begin{equation}
 n_j^{\min,\rm eff}=\max\!\left\{2,
 \min\!\left[n_j^{\min},\left\lceil f_j i_{\rm mobile}\right\rceil\right]\right\}. \label{eq:njeff}
\end{equation}
Equation~\eqref{eq:njeff} is a numerical admissibility rule: it prevents a
calibrated junction channel from becoming identically inaccessible when
$i_{\rm mobile}<n_j^{\min}$.  It is not derived from Marian atomistics and should
be reported as a model regularization.  Combining Eqs.~\eqref{eq:Kii},
\eqref{eq:Kjunc}, and \eqref{eq:phi} gives the requested second kernel.

\subsection{Step 5: mass-action equations and factors of two}

For a symmetric kernel, the product gain at size $p$ is
\begin{equation}
 \left.\dot c_p^{100}\right|_{\rm junc}
 =\frac12\sum_{n=1}^{p-1}K_{n,p-n}^{111\rightarrow100}
 c_n^{111}c_{p-n}^{111},                                               \label{eq:jgain}
\end{equation}
where $1/2$ prevents double counting $(n,p-n)$ and $(p-n,n)$.  The reactant loss is
\begin{equation}
 \left.\dot c_n^{111}\right|_{\rm junc}
 =-c_n^{111}\sum_{n'}K_{nn'}^{111\rightarrow100}c_{n'}^{111}.         \label{eq:jloss}
\end{equation}
Then, after relabeling the double sums,
\begin{equation}
 \left.\frac{\dd}{\dd t}
 \left(\sum_n nc_n^{111}+\sum_p pc_p^{100}\right)\right|_{\rm junc}=0. \label{eq:jcons}
\end{equation}
Loop number is not conserved: every junction converts two loops into one.

\section{Kernel 3: $100_m+111_n\rightarrow100_{m+n}$ absorption}
\label{sec:absorption}

\subsection{Step 1: relative diffusivity and encounter kernel}

The general cross-character Smoluchowski kernel in atomic-fraction units is
\begin{equation}
 \K_{mn}^{100,111}=\frac{8\pi}{\Omega^{2/3}}
 (\xi_m+\xi_n)(D_m^{100}+D_n^{111}).                                  \label{eq:Kcross}
\end{equation}
Because the host $100$ loop is declared sessile,
\begin{equation}
 D_m^{100}=0\quad\Longrightarrow\quad
 \K_{mn}^{100,111}=\frac{8\pi}{\Omega^{2/3}}
 (\xi_m+\xi_n)D_n^{111}.                                             \label{eq:KcrossSessile}
\end{equation}
The asymmetry therefore arises from mobility and product identity, not from
discarding the host size: $m$ remains in the capture radius.

\subsection{Step 2: a distinct templated-completion gate}

The captured $111$ cluster must be incorporated into the character of the much
larger $100$ host.  The same competing-hazard argument as
Eqs.~\eqref{eq:hazardInt}--\eqref{eq:Pgeneral}, but with an absorption-specific
forward barrier and prefactor, gives the closure
\begin{equation}
 \boxed{
 P_{\rm succ}^{\rm abs}(T)=A_{\rm abs}
 \frac{\e^{-\Delta H_2^{\rm abs}/\kb T}}
 {\e^{-\Delta H_2^{\rm abs}/\kb T}+\e^{-\Delta H_{\rm rev}/\kb T}}
 =\frac{A_{\rm abs}}
 {1+\e^{(\Delta H_2^{\rm abs}-\Delta H_{\rm rev})/\kb T}}.}          \label{eq:Pabs}
\end{equation}
This is not a literal probability unless parameters ensure
$0\le P_{\rm succ}^{\rm abs}\le1$.  With $\Delta H_2^{\rm abs}>Delta H_{\rm rev}$,
the high-temperature limit is $A_{\rm abs}/2$, so $A_{\rm abs}\le2$ preserves
the probability interpretation.  More generally, Eq.~\eqref{eq:Pabs} should be
called an \emph{effective success multiplier} or clipped at unity.

The physical reason for allowing $\Delta H_2^{\rm abs}\ne\Delta H_2$ is that
junction formation creates a new $100$ nucleus from two comparable $111$ loops,
whereas absorption is templated by an existing $100$ loop.  This distinction is
a \rc{} hypothesis motivated by the Marian mechanism, not a barrier directly
reported for this exact coarse-grained reaction.

\subsection{Step 3: final absorption kernel and admissibility}

Multiplying Eqs.~\eqref{eq:KcrossSessile} and \eqref{eq:Pabs} gives
\begin{equation}
 \boxed{
 K_{mn}^{\rm abs}=P_{\rm succ}^{\rm abs}(T)
 \frac{8\pi}{\Omega^{2/3}}(\xi_m+\xi_n)D_n^{111},}                    \label{eq:Kabs}
\end{equation}
subject to
\begin{equation}
 m\ge n_{\rm loop}^{\min},\qquad n\le i_{\rm mobile},\qquad m+n\le I. \label{eq:absAdmiss}
\end{equation}
If the implementation contains a separate multiplier
\code{absorb\_boost\_100}, it multiplies the right-hand side of
Eq.~\eqref{eq:Kabs} and should be listed independently from
$P_{\rm succ}^{\rm abs}$.

\subsection{Step 4: master-equation terms and conservation}

The mobile-partner loss is
\begin{equation}
 \left.\dot c_n^{111}\right|_{\rm abs}
 =-c_n^{111}\sum_mK_{mn}^{\rm abs}c_m^{100}.                          \label{eq:abs111loss}
\end{equation}
For the $100$ ladder, absorption moves a loop from $m$ to $m+n$:
\begin{equation}
 \left.\dot c_p^{100}\right|_{\rm abs}
 =\sum_{n=1}^{p-n_{\rm loop}^{\min}}
 K_{p-n,n}^{\rm abs}c_{p-n}^{100}c_n^{111}
 -c_p^{100}\sum_nK_{pn}^{\rm abs}c_n^{111}.                          \label{eq:abs100}
\end{equation}
Multiplying Eqs.~\eqref{eq:abs111loss} and \eqref{eq:abs100} by cluster size and
summing shows
\begin{equation}
 \left.\frac{\dd}{\dd t}
 \left(\sum_nnc_n^{111}+\sum_mmc_m^{100}\right)\right|_{\rm abs}=0. \label{eq:absCons}
\end{equation}
The bottom-boundary rule---a $100$ loop shrinking below
$n_{\rm loop}^{\min}$ returns to the $111$ population---is required to retain
this conservation when other growth, shrinkage, and emission edges operate.

\section{Numerical evaluation of the two success gates}
\label{sec:numerics}

With $\Delta H_2=0.44$~eV, $\Delta H_{\rm rev}=0.30$~eV,
$\Delta H_2^{\rm abs}=0.36$~eV, $A_{\rm abs}=2$, and
$\kb=8.617333262\times10^{-5}$~eV~K$^{-1}$,
Eqs.~\eqref{eq:Psucc} and \eqref{eq:Pabs} give Table~\ref{tab:gates}.

\begin{table}[ht]
\centering
\caption{Recalculated completion factors for the supplied calibrated barriers.}
\label{tab:gates}
\begin{tabular}{rrrr}
\toprule
$T$ ($^\circ$C) & $T$ (K) & $P_{\rm succ}$ & $P_{\rm succ}^{\rm abs}$ \\
\midrule
300 & 573.15 & 0.0555 & 0.4577 \\
330 & 603.15 & 0.0634 & 0.4794 \\
400 & 673.15 & 0.0822 & 0.5245 \\
\bottomrule
\end{tabular}
\end{table}

The absorption multiplier is 8.25, 7.57, and 6.38 times the junction completion
probability at 300, 330, and 400~$^\circ$C.  The ratio decreases with temperature,
although both gates increase.  The actual junction yield is smaller still by
$\varphi_{\max}g(n,n')$.

\section{Complete parameter provenance and consistency table}
\label{sec:data}

\begin{table}[p]
\centering\small
\caption{Inputs to Eqs.~\eqref{eq:Gamma}, \eqref{eq:Kjunc}, and
\eqref{eq:Kabs}.  ``Model/calibration'' means that the numerical value is not
directly supplied by the cited atomistic or elastic theory.}
\label{tab:params}
\begin{tabular}{p{3.6cm}p{2.5cm}p{2.0cm}p{6.4cm}}
\toprule
Quantity & Symbol & Value & Basis/provenance \\
\midrule
Attempt frequency & $\nu_0$ & $10^{13}$ s$^{-1}$ & phonon-scale prefactor; model input \\
Unary barrier intercept & $E_{a0}$ & 2.10 eV & supplied calibration; differs from repository default 1.8 eV \\
Barrier per line segment & $\gamma_a$ & 0.020 eV & supplied calibration; differs from repository default 0.03 eV \\
Loop onset & $n_{\rm loop}^{\min}$ & 4 & population-definition/model input \\
Peak junction efficiency & $\varphi_{\max}$ & 0.10 & unresolved orientation/geometry; calibration \\
Log-size width & $\sigma_s$ & 0.35 & size-compatibility closure; calibration \\
Nominal junction floor & $n_j^{\min}$ & 30 & model/calibration \\
Mobility fraction & $f_j$ & 0.60 & numerical admissibility regularization \\
Junction completion barrier & $\Delta H_2$ & 0.44 eV & calibrated effective barrier inspired by the Marian intermediate \\
Reverse barrier & $\Delta H_{\rm rev}$ & 0.30 eV & calibrated/effective Marian-pathway barrier \\
Absorption completion barrier & $\Delta H_2^{\rm abs}$ & 0.36 eV & distinct templated-absorption closure \\
Absorption prefactor & $A_{\rm abs}$ & 2.0 & calibration; maximum preserving probability interpretation \\
Lattice parameter & $a$ & 2.867 \AA & $\alpha$-Fe material data \\
Atomic volume & $\Omega$ & 11.80 \AA$^3$ & $\alpha$-Fe material data; consistent with $a^3/2$ \\
Core cutoff & $\delta$ & 4.0 \AA & dislocation-core model; \citet{DudarevBullough2008} \\
Softening temperature & $T_c$ & 1185 K & bcc--fcc transition/elastic softening; \citet{DudarevBullough2008} \\
$111$ prelog factor & $\widehat F_{111}$ & 0.28 eV/\AA & digitized/selected from anisotropic-elastic data \\
$100$ zero-$T$ factor & $\widehat F_{100}^{0}$ & 0.5182 eV/\AA & derived from Eq.~\eqref{eq:F100cal}, not independent \\
$111$ core terms & $F_\delta^{111},F_c^{111}$ & 0.345, 0.46 eV/\AA & \citet{DudarevBullough2008} \\
$100$ core terms & $F_\delta^{100},F_c^{100}$ & 0.387, 0.33 eV/\AA & \citet{DudarevBullough2008} \\
Crossover temperature & $T^*$ & 450 $^\circ$C & calibration condition in Eq.~\eqref{eq:F100cal} \\
Reference size & $n_{\rm ref}$ & 30 & calibration condition; supplied analysis value \\
\bottomrule
\end{tabular}
\end{table}

\section{Compact final forms and interpretation}
\label{sec:summary}

The three derived kernels are
\begin{align}
 \Gamma_{\rm uni}(n,T)
 &=\nu_0\e^{-[E_{a0}+\gamma_aP_{111}(n)/b_{111}]/\kb T}
 \pos{1-\e^{-\Delta F(n,T)/\kb T}}
 \Theta(n-n_{\rm loop}^{\min}),                                      \label{eq:final1}\\
 K_{nn'}^{\rm junc}
 &=P_{\rm succ}(T)\varphi_{\max}
 \e^{-\ln^2(n/n')/(2\sigma_s^2)}
 \Theta[\min(n,n')-n_j^{\min,\rm eff}]
 \frac{8\pi}{\Omega^{2/3}}(\xi_n+\xi_{n'})(D_n^{111}+D_{n'}^{111}),  \label{eq:final2}\\
 K_{mn}^{\rm abs}
 &=P_{\rm succ}^{\rm abs}(T)
 \frac{8\pi}{\Omega^{2/3}}(\xi_m+\xi_n)D_n^{111}.                   \label{eq:final3}
\end{align}

Their physical roles are distinct:
\begin{itemize}
\item Equation~\eqref{eq:final1} transfers an existing loop between character
states.  Dudarev theory supplies $\Delta F$; the barrier law and one-way net-rate
interpretation are mesoscale closures.
\item Equation~\eqref{eq:final2} creates a new $100$ loop through an encounter
between two mobile $111$ loops.  Marian atomistics motivates the reaction path;
the log-Gaussian compatibility function and effective barriers coarse-grain
unresolved collision variables.
\item Equation~\eqref{eq:final3} grows an existing sessile $100$ loop by capture
of a mobile $111$ loop.  Its separate completion gate represents templating by
the host and should not be conflated with junction nucleation.
\end{itemize}
All three conserve SIA content by stoichiometry.  The unary channel conserves
loop number; the two binary channels reduce loop number by one per event.  The
clear separation of thermodynamic input, encounter kinetics, and calibrated
branching factors is essential for mechanism-ablation studies and uncertainty
quantification.

\begin{thebibliography}{99}

\bibitem[Dudarev et~al.(2008)Dudarev, Bullough, and Derlet]{DudarevBullough2008}
S.~L. Dudarev, R. Bullough, and P.~M. Derlet,
``Effect of the $\alpha$--$\gamma$ Phase Transition on the Stability of
Dislocation Loops in bcc Iron,'' \emph{Physical Review Letters} \textbf{100},
135503 (2008).
\href{https://doi.org/10.1103/PhysRevLett.100.135503}{doi:10.1103/PhysRevLett.100.135503}.

\bibitem[Marian et~al.(2002)Marian, Wirth, and Perlado]{MarianWirth2002}
J. Marian, B.~D. Wirth, and J.~M. Perlado,
``Mechanism of Formation and Growth of $\langle100\rangle$ Interstitial Loops
in Ferritic Materials,'' \emph{Physical Review Letters} \textbf{88}, 255507
(2002).
\href{https://doi.org/10.1103/PhysRevLett.88.255507}{doi:10.1103/PhysRevLett.88.255507}.

\bibitem[Huang and Marian(2019)]{HuangMarian2019}
S. Huang and J. Marian,
``Rates of diffusion controlled reactions for one-dimensionally-moving species
in 3D space,'' \emph{Philosophical Magazine} \textbf{99}, 2562--2583 (2019).
\href{https://doi.org/10.1080/14786435.2019.1635721}{doi:10.1080/14786435.2019.1635721}.

\bibitem[Kohnert and Wirth(2015)]{KohnertWirth2015}
A.~A. Kohnert and B.~D. Wirth,
``Cluster dynamics models of irradiation damage accumulation in ferritic iron.
II. Effects of reaction dimensionality,'' \emph{Journal of Applied Physics}
\textbf{117}, 154306 (2015).
\href{https://doi.org/10.1063/1.4918316}{doi:10.1063/1.4918316}.

\bibitem[Gao et~al.(2022)Gao, Gaganidze, and Aktaa]{Gao2022}
J. Gao, E. Gaganidze, and J. Aktaa,
``Relative population of $\half\langle111\rangle$ and $\langle100\rangle$
interstitial loops in $\alpha$-Fe under irradiation: Effects of C15 cluster
stability and loop one-dimensional movement,'' \emph{Acta Materialia}
\textbf{233}, 117983 (2022).
\href{https://doi.org/10.1016/j.actamat.2022.117983}{doi:10.1016/j.actamat.2022.117983}.

\bibitem[Smoluchowski(1917)]{Smoluchowski1917}
M. von Smoluchowski,
``Versuch einer mathematischen Theorie der Koagulationskinetik kolloider
Lösungen,'' \emph{Zeitschrift für Physikalische Chemie} \textbf{92}, 129--168
(1917).

\bibitem[Ghoniem(2026)]{RadClusterRepo}
Ghoniem, \emph{RadCluster}, source files
\code{reaction\_rates.py} and \code{eurofer97/declaration.py}, accessed
15 September 2026.
\href{https://github.com/Ghoniem/RadCluster}{github.com/Ghoniem/RadCluster}.

\end{thebibliography}

\end{document}
